\section{Exact elimination and its induced inner product}\label{sec:schur}

The estimates in the preceding sections concern genuine finite-lattice
operators. To use such estimates along a sequence of scales one also needs
a comparison between the fine and coarse excitation problems. This section
proves an abstract comparison; identifying a Yang--Mills renormalization
scheme satisfying all its hypotheses is a separate problem.

\subsection{A one-step comparison}

Let $\mathcal H=P\oplus Q$ be a finite-dimensional complex Hilbert space.
Suppose $H\ge0$ has nonzero kernel and block decomposition
\begin{equation}
 H=\begin{pmatrix}A&B^*\\B&D\end{pmatrix},
 \qquad D\ge dI_Q,\quad d>0 .
\end{equation}
Energies have already been shifted by the \emph{actual} ground energy.
Define
\begin{equation}\label{eq:schur-def}
 C=D^{-1}B,\quad S=A-B^*D^{-1}B,\quad
 Jp=(p,-Cp),\quad G=J^*J=I+C^*C,\quad
 K=G^{-1/2}SG^{-1/2}.
\end{equation}
$S$ is a quadratic energy on the retained coordinates, while $G$ is the
squared norm of their energy-minimizing graph lift. The normalized map
$U=JG^{-1/2}$ is an isometry.

\begin{proposition}[Metric-normalized Schur gap comparison]\label{prop:schur}
If $\operatorname{spec}K\subset\{0\}\cup[\gamma,\infty)$ for $\gamma>0$,
then
\begin{equation}\label{eq:schur-gap}
 H\ge\left(\gamma^{-1}+d^{-1}\right)^{-1}
                (I-\Pi_{\ker H}).
\end{equation}
The kernels of $K$ and $H$ have the same dimension. If $K$ has a nonzero
excited sector, then also $\gap H\le\gap K$.
\end{proposition}

\begin{proof}
For $r=q+Cp$, completion of squares gives
\begin{equation}\label{eq:completion}
 \langle(p,q),H(p,q)\rangle=\langle p,Sp\rangle+\langle r,Dr\rangle.
\end{equation}
Both summands are nonnegative, and hence $\ker H=J\ker S$.
Decompose $p=p_0+p_\perp$ orthogonally in the $G$ inner product,
with $p_0\in\ker S$. Since $\ker K=G^{1/2}\ker S$,
$Jp_\perp\perp\ker H$ and
\[
 \|Jp_\perp\|^2\le \gamma^{-1}\langle p,Sp\rangle.
\]
The decomposition $(p,q)=Jp_0+Jp_\perp+(0,r)$ therefore implies
\begin{align*}
 \operatorname{dist}((p,q),\ker H)
 &\le \|Jp_\perp\|+\|r\|\\
 &\le \gamma^{-1/2}\langle p,Sp\rangle^{1/2}
       +d^{-1/2}\langle r,Dr\rangle^{1/2}.
\end{align*}
Squaring and applying scalar Cauchy--Schwarz proves
\eqref{eq:schur-gap}. The isometry $U$ identifies $\ker K$ with $\ker H$.
For the upper comparison, minimize the fine Rayleigh quotient over
$U((\ker K)^\perp)$; on that subspace it equals the coarse Rayleigh quotient.
\end{proof}

\begin{remark}[Closed-form extension]\label{rem:schur-domain}
The same proof applies to a closed nonnegative form on $P\oplus Q$
when $D\ge dI$ is self-adjoint, $C:P\to Q$ is bounded, $s$ is a closed
nonnegative form, and the fine form has exactly the representation
\[
 h[(p,q)]=s[p]+\|D^{1/2}(q+Cp)\|^2
\]
on the domain
\[
 \{(p,q):p\in\Dom s,\ q+Cp\in\Dom D^{1/2}\}.
\]
Here $\Dom s$ denotes the form domain. The coarse operator is the
self-adjoint operator associated with
$k[z]=s[G^{-1/2}z]$ on $G^{1/2}\Dom s$, not an unchecked product
of unbounded operators. The distance argument extends on this domain.
An application to an unbounded lattice Hamiltonian must verify these
conditions; a finite-matrix computation does not verify them.
No extension to unbounded $C$ is asserted.
\end{remark}

\subsection{Why the metric matters}

For $M\in\mathbb R$ consider
\begin{equation}
 H_M=\begin{pmatrix}
 0&0&0\\
 0&1+M^2&M\\
 0&M&1
 \end{pmatrix},\qquad
 P=\operatorname{span}\{e_1,e_2\},\quad Q=\operatorname{span}\{e_3\}.
\end{equation}
Then $D=1$, $S=\operatorname{diag}(0,1)$ and
$G=\operatorname{diag}(1,1+M^2)$, whereas
\begin{equation}
 \gap H_M=\frac{M^2+2-\sqrt{M^4+4M^2}}2\longrightarrow0
 \quad\text{as }|M|\longrightarrow\infty .
\end{equation}
Thus positive lower bounds for $D$ and the \emph{unweighted} Schur
operator $S$ give no positive uniform lower bound for the fine gap.
The correct coarse gap is $(1+M^2)^{-1}$, and
Proposition~\ref{prop:schur} yields the consistent lower bound
$(M^2+2)^{-1}$.

\subsection{An additive inverse-gap budget}

\begin{corollary}[Conditional multiscale bound]\label{cor:multiscale}
Suppose $H_0,\ldots,H_k$ form a sequence of exact metric-normalized
Schur reductions satisfying Proposition~\ref{prop:schur}, or its
closed-form extension, in the same physical energy units. At step $j$
assume the discarded block satisfies $D_j\ge d_jI$ with $d_j>0$.
If $H_k$ has gap at least $m_*>0$, then
\begin{equation}\label{eq:budget}
 \gap H_0\ge
 \left(m_*^{-1}+\sum_{j=0}^{k-1}d_j^{-1}\right)^{-1}.
\end{equation}
\end{corollary}
\begin{proof}
Iterate \eqref{eq:schur-gap} backwards from the terminal operator.
\end{proof}

This avoids a fixed multiplicative loss per scale, but every operator and
metric in the hypothesis must be the actual one generated by elimination.
For example, assume \emph{as a model profile}
$a_r=AL^{-r}$, $x_r=(\sigma+qr)^{-1}$ for $r\ge1$,
where $A,\sigma>0$, $L>1$ and $q\ge0$. If one could prove
$d_r\ge cx_r/a_r$ with one $c>0$, then
\begin{equation}
 \sum_{r\ge1}d_r^{-1}
 \le\frac{A}{c}\left(\frac{\sigma}{L-1}
                   +\frac{qL}{(L-1)^2}\right).
\end{equation}
The geometric summation is exact. Neither the assumed trajectory nor
the nonlinear discarded-sector estimate is established here.

\subsection{Quantum Schur elimination need not preserve quadratic form}

A separate obstruction is that the normalized reduction above is not
the same operation as integrating a Gaussian coordinate density.
Let $H=\mathrm d\Gamma(h)$ be the bosonic second quantization of
\[
 h=\begin{pmatrix}2&1\\1&2\end{pmatrix}.
\]
Retain the subspace spanned by all $|n,0\rangle$, so that the retained
space contains the true Fock vacuum. Total particle number is preserved
by $H$ and by all the reduction maps. In the one-particle sector,
\[
 S_1=\frac32,\qquad G_1=\frac54,\qquad K_1=\frac65 .
\]
In the two-particle sector the fine matrix is
\[
 H_2=\begin{pmatrix}
 4&\sqrt2&0\\
 \sqrt2&4&\sqrt2\\
 0&\sqrt2&4
 \end{pmatrix}
\]
in the basis $|2,0\rangle,|1,1\rangle,|0,2\rangle$. Direct inversion gives
\[
 S_2=\frac{24}{7},\qquad G_2=\frac{58}{49},
 \qquad K_2=\frac{84}{29}\ne 2K_1.
\]
Since $h\ge I$, the normalized graph restriction in the $n$-particle
sector is bounded below by $nI$. Consequently $K_1$ and $K_2<3$ are
indeed the first two positive levels. They are not equally spaced.
The full Fock-space form construction is legitimate: on the discarded
part of the $n$-particle sector, $D_n\ge nI$ and $\|B_n\|=\sqrt n$,
so $\|D_n^{-1}B_n\|\le n^{-1/2}$ and the graph operator is bounded.
This example rules out automatic closure of the oscillator class under
this quantum elimination. It does not rule out controlling the generated
operators in a larger class.
